% Generated draft: requires mathematical and source review.
\documentclass[11pt]{article}
\usepackage[margin=25mm]{geometry}
\usepackage{fontspec}
\setmainfont{Latin Modern Roman}
\usepackage{amsmath,amssymb,mathtools,unicode-math}
\setmathfont{Latin Modern Math}
\usepackage{enumitem,needspace,xurl,hyperref}
\setlength{\parindent}{0pt}
\setlength{\parskip}{5pt}
\setlength{\emergencystretch}{3em}
\newcommand{\N}{\mathbb{N}}
\newcommand{\Z}{\mathbb{Z}}
\newcommand{\Q}{\mathbb{Q}}
\newcommand{\R}{\mathbb{R}}
\newcommand{\C}{\mathbb{C}}
\newcommand{\F}{\mathbb{F}}
\newcommand{\A}{\mathbb{A}}
\newcommand{\PP}{\mathbb{P}}
\newcommand{\E}{\mathbb{E}}
\newcommand{\eps}{\varepsilon}
\newcommand{\dd}{\,\mathrm{d}}
\begin{document}
\section*{UnsolvedMath ID: 20000221}
\textbf{Source problem code:} catalogue entry
\subsection{Definitions and mathematical statement}
Brian Osserman asks whether a curve $X$ formed by gluing two rational curves at $g+1$ generic pairs of points is Brill--Noether general. Concretely, take $C_1\simeq C_2\simeq\PP^1$ with ordered pairwise distinct points $p_0,\ldots,p_g$ and $q_0,\ldots,q_g$, and identify $p_i\sim q_i$ transversely. The resulting connected binary nodal curve has arithmetic genus $g$. ``Generic'' must mean independently general labeled configurations on the two components, rather than configurations related by one projectivity.

For a multidegree $(a,b)$, let $G^r_{(a,b)}(X)$ parametrize line bundles $L$ with component degrees $a,b$ together with $(r+1)$-dimensional subspaces $V\subset H^0(X,L)$, and put
\[
d=a+b,\qquad \rho(g,r,d)=g-(r+1)(g-d+r).
\]
An expected-dimension interpretation asks that the chosen linear-series spaces have dimension $\rho$ when nonempty and be empty for $\rho<0$. On a reducible curve one must specify which multidegrees, balancedness condition, compactification or limit-linear-series theory, and whether nonemptiness for $\rho\ge0$ is part of the assertion. The source does not choose these conventions or a field and genus range.

The attached unreviewed aid attributes an affirmative intended expected-dimension answer for general binary curves to Xiang He's all-rank theorem, including every potentially nonempty balanced multidegree and Osserman limit linear series. Its exact hypotheses and full reference are not supplied. It cautions that the unqualified question is definition-dependent.

Its separate boundary calculation assumes $a<0$. For every binary curve it reports $G^r_{(a,b)}(X)=\varnothing$ if $b<g+r+1$; otherwise it is a Grassmann bundle over the multidegree Picard torsor, with
\[
\dim G^r_{(a,b)}=g+(r+1)(b-g-r-1),\qquad
\dim G^r_{(a,b)}-\rho=(r+1)(-a-1).
\]
It reports balanced examples at $a=-1$ with $\rho\ge0$ but empty $G^0$, and positive dimension excess below $a=-1$ when nonempty. Finally it claims that a $(1,1)$-pencil exists exactly when the two labeled node configurations are projectively equivalent; for $g\ge2$ this diagonal in their configuration moduli has codimension $g-2$. The genus restriction is necessary for that codimension expression, and the pencil convention must be fixed. These diagnostics do not justify an unrestricted all-multidegree generality assertion.
\subsection{Short English statement}
Ask whether independently general gluing of two rational components produces a Brill--Noether general binary curve. The answer depends on the chosen multidegrees and linear-series theory; the aid reports an affirmative theorem in the intended balanced setting and explicit boundary exceptions to stronger readings.
\subsection{Sources}
\begin{itemize}
\item Brian Osserman, suggested Problem 4 in AIM \emph{Degenerations in algebraic geometry}: \url{https://aimath.org/pastworkshops/degenalggeomproblems.pdf}. Canonical location: \texttt{aim-algebraic-geometry-notes.json}, notes[220]. Original wording is labeled exact; remarks and literature are empty. The supplied extraction is used without consulting a PDF.
\item The attached aid cites Xiang He's all-rank theorem and Osserman limit linear series without complete bibliographic entries, and proposes the negative-multidegree and node-configuration diagnostics with low novelty confidence. It is machine generated and not peer reviewed; the source's solved-in-literature label is retained as metadata.
\item Local registry: \texttt{data/unsolvedmath/problems.json}, permanent catalogue entry, exact source code \texttt{catalogue entry}.
\end{itemize}
\end{document}
